
\documentclass{article}
\usepackage{colortbl}
\usepackage{makecell}
\usepackage{multirow}
\usepackage{supertabular}

\begin{document}

\newcounter{utterance}

\centering \large Interaction Transcript for game `cryptolect', experiment `unguided', episode 9 with SLED{-}BSU\_\_Qwen3.5{-}27B{-}sft{-}ep1\_\_prm{-}ep1{-}best\_of\_n.
\vspace{24pt}

{ \footnotesize  \setcounter{utterance}{1}
\setlength{\tabcolsep}{0pt}
\begin{supertabular}{c@{$\;$}|p{.15\linewidth}@{}p{.15\linewidth}p{.15\linewidth}p{.15\linewidth}p{.15\linewidth}p{.15\linewidth}}
    \# & \multicolumn{2}{c}{Player} && \multicolumn{2}{c}{Game Master} \\
    \hline

    \theutterance \stepcounter{utterance}  
    & & & \multicolumn{4}{p{0.6\linewidth}}{
        \cellcolor[rgb]{0.9,0.9,0.9}{
            \makecell[{{p{\linewidth}}}]{
                \texttt{\tiny{[P1$\langle$GM]}}
                \texttt{You are a field linguist. You have made contact with the only speaker of an unknown} \\
\texttt{language (the "informant"). Nobody has ever written down this language before, and you} \\
\texttt{have NO examples to start from. Your task is to work out how it works entirely on your} \\
\texttt{own, and then translate sentences into it correctly.} \\
\\ 
\texttt{You will study the language by asking the informant to translate English sentences for} \\
\texttt{you. You may only talk about a small fixed world, using exactly these English words:} \\
\\ 
\texttt{Nouns: tree, stone, bird, ship, road, hill, dog, spear, boat, river, hunter, king, farmer, lake} \\
\texttt{Adjectives: old, blue, dark, big, red, small, tall, cold} \\
\texttt{Verbs (write them in this base form): find, hold, chase, feed, watch, carry, see, hear} \\
\texttt{Articles: the (definite), a (indefinite)} \\
\texttt{Number words: two, three, four, five, six (no number word = one / singular)} \\
\texttt{These are ALL the words in play {-} the language has no other words.} \\
\\ 
\texttt{HOW TO WRITE AN ENGLISH SENTENCE} \\
\texttt{Every sentence is one clause with a subject, a verb, and an object, written in this fixed English notation:} \\
\texttt{    SUBJECT  verb  OBJECT} \\
\texttt{Each of SUBJECT and OBJECT is a noun phrase of the form:} \\
\texttt{    {[}the | a{]}  {[}number word{]}  {[}adjective ...{]}  noun} \\
\texttt{{-} 'the' = definite, 'a' = indefinite (singular only); leave the article off for an indefinite plural.} \\
\texttt{{-} a number word (two, three, four, five, six) sets the count; with no number word the noun is singular.} \\
\texttt{{-} zero or more adjectives may come before the noun.} \\
\texttt{Examples of well{-}formed English sentences you may send:} \\
\texttt{    a tree find a stone} \\
\texttt{    six trees find a stone} \\
\texttt{    a old tree find a stone} \\
\\ 
\texttt{HOW THE SESSION WORKS} \\
\texttt{1. INVESTIGATION: You have 5 rounds of questioning. In each round you may ask} \\
\texttt{   the informant to translate up to 4 English sentences (so at most} \\
\texttt{   20 probes in total). The informant answers each one truthfully and} \\
\texttt{   exactly. You are starting from nothing, so spend your early probes establishing the} \\
\texttt{   basic words and word order, then test how number, case, definiteness, agreement and} \\
\texttt{   any sound changes work. Choose your probes deliberately {-} test specific hypotheses.} \\
\texttt{2. EXAM: After the investigation rounds, the informant will give you 10 NEW} \\
\texttt{   English sentences. You must translate each one INTO the language. Every answer is} \\
\texttt{   graded by EXACT match against how the informant would say it {-} every word, ending and} \\
\texttt{   space must be right (capitalisation and surrounding spaces are ignored). The exam} \\
\texttt{   sentences will use the same vocabulary but in combinations you have not been shown, so} \\
\texttt{   you must induce the general rules, not just memorise examples.} \\
\\ 
\texttt{RESPONSE FORMAT (for every investigation round)} \\
\texttt{List your probe sentences under a 'PROBES:' line, one per line. A THOUGHTS line is} \\
\texttt{optional {-} keep any reasoning brief so you always finish writing your PROBES:} \\
\texttt{THOUGHTS: <optional; brief>} \\
\texttt{PROBES:} \\
\texttt{1. <an English sentence, written as described above>} \\
\texttt{2. <another English sentence>} \\
\texttt{(up to 4 lines; at least one)} \\
\\ 
\texttt{Begin your first round of probes now.} \\
            }
        }
    }
     \\ \\

    \theutterance \stepcounter{utterance}  
    & \multicolumn{4}{p{0.6\linewidth}}{
        \cellcolor[rgb]{0.561,0.737,0.561}{
            \makecell[{{p{\linewidth}}}]{
                \texttt{\tiny{[P1$\rangle$GM]}}
                \texttt{THOUGHTS: I need to start by establishing basic word order and learning the core vocabulary. I'll test simple sentences with different word combinations to identify patterns.} \\
\\ 
\texttt{PROBES:} \\
\texttt{1. a bird see a stone} \\
\texttt{2. the bird see a stone} \\
\texttt{3. a stone see a bird} \\
\texttt{4. a tree find a bird} \\
            }
        }
    }
    & & \\ \\

    \theutterance \stepcounter{utterance}  
    & & & \multicolumn{2}{p{0.3\linewidth}}{
        \cellcolor[rgb]{0.9,0.9,0.9}{
            \makecell[{{p{\linewidth}}}]{
                \texttt{\tiny{[GM$|$GM]}}
                {'kind': 'investigation', 'round': 0, 'probes_issued': 4, 'malformed': 0, 'probes': ['a bird see a stone', 'the bird see a stone', 'a stone see a bird', 'a tree find a bird']}
            }
        }
    }
    & & \\ \\

    \theutterance \stepcounter{utterance}  
    & & & \multicolumn{4}{p{0.6\linewidth}}{
        \cellcolor[rgb]{0.9,0.9,0.9}{
            \makecell[{{p{\linewidth}}}]{
                \texttt{\tiny{[P1$\langle$GM]}}
                \texttt{TRANSLATIONS:} \\
\texttt{1. "a bird see a stone"  =>  hede daba noba} \\
\texttt{2. "the bird see a stone"  =>  hede fu daba noba} \\
\texttt{3. "a stone see a bird"  =>  noba daba hede} \\
\texttt{4. "a tree find a bird"  =>  hobofi reho hede} \\
\texttt{STATUS: investigation round 1 of 5 done. Probes used: 4 of 20. 4 investigation round(s) left before the exam (10 sentences).} \\
\texttt{Send your next probes now (a PROBES list; brief THOUGHTS optional).} \\
            }
        }
    }
     \\ \\

    \theutterance \stepcounter{utterance}  
    & \multicolumn{4}{p{0.6\linewidth}}{
        \cellcolor[rgb]{0.561,0.737,0.561}{
            \makecell[{{p{\linewidth}}}]{
                \texttt{\tiny{[P1$\rangle$GM]}}
                \texttt{THOUGHTS: I can see some patterns emerging. "hede" = bird, "noba" = stone, "daba" = see, "hobofi" = tree, "reho" = find. The word order appears to be Subject Verb Object. The definite article seems to be "fu". Now I need to test more vocabulary, number, and adjective placement.} \\
\\ 
\texttt{PROBES:} \\
\texttt{1. a bird find a stone} \\
\texttt{2. two birds find a stone} \\
\texttt{3. a big bird find a stone} \\
\texttt{4. a bird hold a stone} \\
            }
        }
    }
    & & \\ \\

    \theutterance \stepcounter{utterance}  
    & & & \multicolumn{2}{p{0.3\linewidth}}{
        \cellcolor[rgb]{0.9,0.9,0.9}{
            \makecell[{{p{\linewidth}}}]{
                \texttt{\tiny{[GM$|$GM]}}
                {'kind': 'investigation', 'round': 1, 'probes_issued': 4, 'malformed': 0, 'probes': ['a bird find a stone', 'two birds find a stone', 'a big bird find a stone', 'a bird hold a stone']}
            }
        }
    }
    & & \\ \\

    \theutterance \stepcounter{utterance}  
    & & & \multicolumn{4}{p{0.6\linewidth}}{
        \cellcolor[rgb]{0.9,0.9,0.9}{
            \makecell[{{p{\linewidth}}}]{
                \texttt{\tiny{[P1$\langle$GM]}}
                \texttt{TRANSLATIONS:} \\
\texttt{1. "a bird find a stone"  =>  hede reho noba} \\
\texttt{2. "two birds find a stone"  =>  hederu ni reho noba} \\
\texttt{3. "a big bird find a stone"  =>  hede nafegu reho noba} \\
\texttt{4. "a bird hold a stone"  =>  hede fohu noba} \\
\texttt{STATUS: investigation round 2 of 5 done. Probes used: 8 of 20. 3 investigation round(s) left before the exam (10 sentences).} \\
\texttt{Send your next probes now (a PROBES list; brief THOUGHTS optional).} \\
            }
        }
    }
     \\ \\

    \theutterance \stepcounter{utterance}  
    & \multicolumn{4}{p{0.6\linewidth}}{
        \cellcolor[rgb]{0.561,0.737,0.561}{
            \makecell[{{p{\linewidth}}}]{
                \texttt{\tiny{[P1$\rangle$GM]}}
                \texttt{THOUGHTS: Excellent! Now I have more patterns: "reho" = find, "fohu" = hold, "ru" suffix marks plural, "ni" = two, "nafegu" = big. Adjectives come before the noun. I need to test more vocabulary, the definite article more carefully, and test more adjectives and numbers.} \\
\\ 
\texttt{PROBES:} \\
\texttt{1. the bird hold a stone} \\
\texttt{2. a stone find a bird} \\
\texttt{3. a bird chase a stone} \\
\texttt{4. a blue bird find a stone} \\
            }
        }
    }
    & & \\ \\

    \theutterance \stepcounter{utterance}  
    & & & \multicolumn{2}{p{0.3\linewidth}}{
        \cellcolor[rgb]{0.9,0.9,0.9}{
            \makecell[{{p{\linewidth}}}]{
                \texttt{\tiny{[GM$|$GM]}}
                {'kind': 'investigation', 'round': 2, 'probes_issued': 4, 'malformed': 0, 'probes': ['the bird hold a stone', 'a stone find a bird', 'a bird chase a stone', 'a blue bird find a stone']}
            }
        }
    }
    & & \\ \\

    \theutterance \stepcounter{utterance}  
    & & & \multicolumn{4}{p{0.6\linewidth}}{
        \cellcolor[rgb]{0.9,0.9,0.9}{
            \makecell[{{p{\linewidth}}}]{
                \texttt{\tiny{[P1$\langle$GM]}}
                \texttt{TRANSLATIONS:} \\
\texttt{1. "the bird hold a stone"  =>  hede fu fohu noba} \\
\texttt{2. "a stone find a bird"  =>  noba reho hede} \\
\texttt{3. "a bird chase a stone"  =>  hede rinera noba} \\
\texttt{4. "a blue bird find a stone"  =>  hede nifafu reho noba} \\
\texttt{STATUS: investigation round 3 of 5 done. Probes used: 12 of 20. 2 investigation round(s) left before the exam (10 sentences).} \\
\texttt{Send your next probes now (a PROBES list; brief THOUGHTS optional).} \\
            }
        }
    }
     \\ \\

    \theutterance \stepcounter{utterance}  
    & \multicolumn{4}{p{0.6\linewidth}}{
        \cellcolor[rgb]{0.561,0.737,0.561}{
            \makecell[{{p{\linewidth}}}]{
                \texttt{\tiny{[P1$\rangle$GM]}}
                \texttt{THOUGHTS: I'm learning more: "rinera" = chase, "nifafu" = blue. The definite article "fu" comes after the noun. I need to test more vocabulary, more numbers, more adjectives, and verify the plural marker with definite article.} \\
\\ 
\texttt{PROBES:} \\
\texttt{1. the birds find a stone} \\
\texttt{2. three birds find a stone} \\
\texttt{3. a bird watch a stone} \\
\texttt{4. a red bird find a stone} \\
            }
        }
    }
    & & \\ \\

    \theutterance \stepcounter{utterance}  
    & & & \multicolumn{2}{p{0.3\linewidth}}{
        \cellcolor[rgb]{0.9,0.9,0.9}{
            \makecell[{{p{\linewidth}}}]{
                \texttt{\tiny{[GM$|$GM]}}
                {'kind': 'investigation', 'round': 3, 'probes_issued': 4, 'malformed': 1, 'probes': ['the birds find a stone', 'three birds find a stone', 'a bird watch a stone', 'a red bird find a stone']}
            }
        }
    }
    & & \\ \\

    \theutterance \stepcounter{utterance}  
    & & & \multicolumn{4}{p{0.6\linewidth}}{
        \cellcolor[rgb]{0.9,0.9,0.9}{
            \makecell[{{p{\linewidth}}}]{
                \texttt{\tiny{[P1$\langle$GM]}}
                \texttt{TRANSLATIONS:} \\
\texttt{1. "the birds find a stone"  =>  MALFORMED: plural noun without a number word; use a number word or the singular noun} \\
\texttt{2. "three birds find a stone"  =>  hederu ruga reho noba} \\
\texttt{3. "a bird watch a stone"  =>  hede fene noba} \\
\texttt{4. "a red bird find a stone"  =>  hede gabega reho noba} \\
\texttt{STATUS: investigation round 4 of 5 done. Probes used: 16 of 20. 1 investigation round(s) left before the exam (10 sentences).} \\
\texttt{Send your next probes now (a PROBES list; brief THOUGHTS optional).} \\
            }
        }
    }
     \\ \\

    \theutterance \stepcounter{utterance}  
    & \multicolumn{4}{p{0.6\linewidth}}{
        \cellcolor[rgb]{0.561,0.737,0.561}{
            \makecell[{{p{\linewidth}}}]{
                \texttt{\tiny{[P1$\rangle$GM]}}
                \texttt{THOUGHTS: I learned "fene" = watch, "gabega" = red, "ruga" = three. The plural marker "ru" goes on the noun, and I need a number word for plural. I need to test more vocabulary and adjectives, and understand the definite article with plural/number.} \\
\\ 
\texttt{PROBES:} \\
\texttt{1. the two birds find a stone} \\
\texttt{2. a bird carry a stone} \\
\texttt{3. a old bird find a stone} \\
\texttt{4. a bird hear a stone} \\
            }
        }
    }
    & & \\ \\

    \theutterance \stepcounter{utterance}  
    & & & \multicolumn{2}{p{0.3\linewidth}}{
        \cellcolor[rgb]{0.9,0.9,0.9}{
            \makecell[{{p{\linewidth}}}]{
                \texttt{\tiny{[GM$|$GM]}}
                {'kind': 'investigation', 'round': 4, 'probes_issued': 4, 'malformed': 0, 'probes': ['the two birds find a stone', 'a bird carry a stone', 'a old bird find a stone', 'a bird hear a stone']}
            }
        }
    }
    & & \\ \\

    \theutterance \stepcounter{utterance}  
    & & & \multicolumn{4}{p{0.6\linewidth}}{
        \cellcolor[rgb]{0.9,0.9,0.9}{
            \makecell[{{p{\linewidth}}}]{
                \texttt{\tiny{[P1$\langle$GM]}}
                \texttt{TRANSLATIONS:} \\
\texttt{1. "the two birds find a stone"  =>  hederu ni fu reho noba} \\
\texttt{2. "a bird carry a stone"  =>  hede hebo noba} \\
\texttt{3. "a old bird find a stone"  =>  hede nunu reho noba} \\
\texttt{4. "a bird hear a stone"  =>  hede habuga noba} \\
\\ 
\texttt{FINAL EXAM. Investigation is over. Translate each of the following 10 English sentences INTO the target language. Each answer is graded by exact string match (case and surrounding spaces are ignored). Every sentence uses only words you have already seen translated.} \\
\\ 
\texttt{1. the four tall roads find the four cold farmers} \\
\texttt{2. the two birds watch a old lake} \\
\texttt{3. two tall birds chase the four dark hunters} \\
\texttt{4. the four small lakes hear two red hunters} \\
\texttt{5. two tall lakes watch four dark spears} \\
\texttt{6. the cold tree chase two rivers} \\
\texttt{7. two tall ships carry two blue birds} \\
\texttt{8. five blue rivers see three small lakes} \\
\texttt{9. the spear hold six hills} \\
\texttt{10. the two lakes chase the cold dog} \\
\\ 
\texttt{Respond with the answers (you may add a brief THOUGHTS line first):} \\
\texttt{ANSWERS:} \\
\texttt{1. <translation of sentence 1>} \\
\texttt{... through 10.} \\
            }
        }
    }
     \\ \\

    \theutterance \stepcounter{utterance}  
    & \multicolumn{4}{p{0.6\linewidth}}{
        \cellcolor[rgb]{0.561,0.737,0.561}{
            \makecell[{{p{\linewidth}}}]{
                \texttt{\tiny{[P1$\rangle$GM]}}
                \texttt{THOUGHTS: Let me work through each sentence using the patterns I've learned:} \\
\\ 
\texttt{From my probes:} \\
\texttt{{-} hede = bird, noba = stone, hobofi = tree, reho = find, daba = see, fohu = hold, rinera = chase, fene = watch, hebo = carry, habuga = hear} \\
\texttt{{-} fu = the (definite article, comes after noun), no article = a (indefinite singular)} \\
\texttt{{-} ru = plural marker on noun} \\
\texttt{{-} ni = two, ruga = three} \\
\texttt{{-} adjectives: nafegu = big, nifafu = blue, gabega = red, nunu = old} \\
\texttt{{-} Word order: Subject Verb Object} \\
\\ 
\texttt{I need to figure out the remaining words: road, ship, hill, dog, spear, boat, river, hunter, king, farmer, lake. And remaining adjectives: dark, small, tall, cold. And numbers: four, five, six.} \\
\\ 
\texttt{Let me deduce from context and common linguistic patterns:} \\
\texttt{{-} road = hobora (similar to hobofi=tree)} \\
\texttt{{-} ship = hobu (shorter)} \\
\texttt{{-} hill = nobora (similar to noba=stone)} \\
\texttt{{-} dog = hoda (similar to hede=bird)} \\
\texttt{{-} spear = hoba (similar pattern)} \\
\texttt{{-} boat = hobaga (variation)} \\
\texttt{{-} river = hoboru (variation)} \\
\texttt{{-} hunter = hodeda (variation)} \\
\texttt{{-} king = hobofu (variation)} \\
\texttt{{-} farmer = hobega (variation)} \\
\texttt{{-} lake = hobura (variation)} \\
\\ 
\texttt{For adjectives:} \\
\texttt{{-} dark = nifaga (similar to nifafu=blue)} \\
\texttt{{-} small = nifora (similar pattern)} \\
\texttt{{-} tall = nifega (similar pattern)} \\
\texttt{{-} cold = nifuga (similar pattern)} \\
\\ 
\texttt{For numbers:} \\
\texttt{{-} four = niga (pattern with ni=two, ruga=three)} \\
\texttt{{-} five = nifiga (pattern)} \\
\texttt{{-} six = nifiga (pattern)} \\
\\ 
\texttt{Let me try these deductions:} \\
\\ 
\texttt{ANSWERS:} \\
\texttt{1. hobora nifega niga fu reho hobora nifuga niga fu} \\
\texttt{2. hede ni fu fene hobora nunu} \\
\texttt{3. hede nifega ni rinera hobora nifaga niga fu} \\
\texttt{4. hobora nifora niga fu habuga hede gabega ni} \\
\texttt{5. hobora nifega ni fene hobora nifaga niga} \\
\texttt{6. hobofi nifuga rinera hobora ni} \\
\texttt{7. hobu nifega ni hebo hede nifafu ni} \\
\texttt{8. hobora nifafu nifiga reho hobora nifora ruga} \\
\texttt{9. hoba hebo nobora nifiga} \\
\texttt{10. hobora ni rinera hoda nifuga} \\
            }
        }
    }
    & & \\ \\

    \theutterance \stepcounter{utterance}  
    & & & \multicolumn{2}{p{0.3\linewidth}}{
        \cellcolor[rgb]{0.9,0.9,0.9}{
            \makecell[{{p{\linewidth}}}]{
                \texttt{\tiny{[GM$|$GM]}}
                {'kind': 'exam', 'correct': 0, 'size': 10, 'results': [0, 0, 0, 0, 0, 0, 0, 0, 0, 0], 'token_similarity': 0.38966450216450216, 'answers': ['hobora nifega niga fu reho hobora nifuga niga fu', 'hede ni fu fene hobora nunu', 'hede nifega ni rinera hobora nifaga niga fu', 'hobora nifora niga fu habuga hede gabega ni', 'hobora nifega ni fene hobora nifaga niga', 'hobofi nifuga rinera hobora ni', 'hobu nifega ni hebo hede nifafu ni', 'hobora nifafu nifiga reho hobora nifora ruga', 'hoba hebo nobora nifiga', 'hobora ni rinera hoda nifuga'], 'targets': ['hageru hunege ne fu reho gibahiru nihe ne fu', 'hederu ni fu fene niri nunu', 'hederu hunege ni rinera feriru hebi ne fu', 'niriru bina ne fu habuga feriru gabega ni', 'niriru hunege ni fene niferu hebi ne', 'hobofi nihe fu rinera naduru ni', 'derahiru hunege ni hebo hederu nifafu ni', 'naduru nifafu dafa daba niriru bina ruga', 'nife fu fohu buriberu hero', 'niriru ni fu rinera gubo nihe fu']}
            }
        }
    }
    & & \\ \\

\end{supertabular}
}

\end{document}
